% GENERATED FILE. Edit canonical lessons, then run npm run generate:books.
% canonical-source-hash: fnv1a64:5d63282d83c7b14e
% canonical-lessons: FR-C493-grossesse, FR-C493-accoucher, FR-C493-nourrisson, FR-C493-biberon, FR-C493-bercer

\chapter{La première année}
\label{ch:fr-la-premiere-annee}
\hlchaptermodality{493}

\begin{chapteropening}
\textbf{By the end of this chapter:} \emph{I can talk about pregnancy and a baby's first months.}

Complete the last lesson of chapter 493: I can talk about pregnancy and a baby's first months.
\end{chapteropening}

\section[la grossesse]{la grossesse — pregnancy}

\label{lesson:FR-C493-grossesse}



\begin{quote}
Before the new one: say the French for siblings, then the French for large family.
\end{quote}

\subsection*{You'll want to know: la grossesse}
\textbf{la grossesse} — \textquotedblleft{}pregnancy\textquotedblright{}.

\begin{quote}
Pendant sa grossesse, elle a continué à travailler.
\end{quote}

\begin{quote}
\emph{During her pregnancy she kept on working.}
\end{quote}

\emph{Être enceinte} — to be pregnant.

\begin{cousinweb}
\emph{la grossesse} is built on \emph{gros} \textquotedblleft{}big\textquotedblright{} with \emph{-esse}; \emph{gros} is inherited from Late Latin \emph{grossus} \textquotedblleft{}thick\textquotedblright{}.
\end{cousinweb}

\begin{grammarlens}[title={What you've built}]
One more word for telling, explaining and describing at B1.
\end{grammarlens}

\subsection*{Your turn}
\begin{itemize}
\raggedright
  \item \emph{Say it:} \emph{la grossesse}
  \item \emph{Say it:} \emph{la grossesse}, once more
  \item \emph{Say it:} a sentence of your own with \emph{la grossesse}
  \item \emph{Recall:} say the French for siblings, then the French for large family, then say \emph{la grossesse} again
\end{itemize}

\begin{tcolorbox}[breakable,colback=teal!4,colframe=teal!35!black,title={Before you move on}]
What does la grossesse mean? (\textquotedblleft{}Pregnancy\textquotedblright{}.) Say it once more.
\end{tcolorbox}
\section[accoucher]{accoucher — to give birth}

\label{lesson:FR-C493-accoucher}



\begin{quote}
Before the new one: say the French for large family, then the French for pregnancy.
\end{quote}

\subsection*{You'll want to know: accoucher}
\textbf{accoucher} — \textquotedblleft{}to give birth\textquotedblright{}.

\begin{quote}
Ma sœur a accouché d'un garçon ce matin.
\end{quote}

\begin{quote}
\emph{My sister gave birth to a boy this morning.}
\end{quote}

\emph{Accoucher de} — to give birth to. \emph{L'accouchement} — the birth.

\begin{cousinweb}
\emph{accoucher} comes from Old French \emph{acouchier} \textquotedblleft{}to lay down, to take to one's bed\textquotedblright{} (EtymDB), built on \emph{coucher}, from Latin \emph{collocāre} \textquotedblleft{}to place\textquotedblright{}, itself built on \emph{locus} \textquotedblleft{}place\textquotedblright{} (Lewis \& Short).
\end{cousinweb}

\begin{grammarlens}[title={What you've built}]
One more word for telling, explaining and describing at B1.
\end{grammarlens}

\subsection*{Your turn}
\begin{itemize}
\raggedright
  \item \emph{Say it:} \emph{accoucher}
  \item \emph{Say it:} \emph{accoucher}, once more
  \item \emph{Say it:} a sentence of your own with \emph{accoucher}
  \item \emph{Recall:} say the French for large family, then the French for pregnancy, then say \emph{accoucher} again
\end{itemize}

\begin{tcolorbox}[breakable,colback=teal!4,colframe=teal!35!black,title={Before you move on}]
What does accoucher mean? (\textquotedblleft{}To give birth\textquotedblright{}.) Say it once more.
\end{tcolorbox}
\section[le nourrisson]{le nourrisson — infant, small baby}

\label{lesson:FR-C493-nourrisson}



\begin{quote}
Before the new one: say the French for pregnancy, then the French for to give birth.
\end{quote}

\subsection*{You'll want to know: le nourrisson}
\textbf{le nourrisson} — \textquotedblleft{}infant, small baby\textquotedblright{}.

\begin{quote}
Le nourrisson dort presque toute la journée.
\end{quote}

\begin{quote}
\emph{The infant sleeps almost all day.}
\end{quote}

Formal: a baby of a few months.

\begin{cousinweb}
\emph{le nourrisson} comes from Late Latin \emph{nūtrītiō} \textquotedblleft{}feeding, nursing\textquotedblright{}, built on \emph{nūtrīre} \textquotedblleft{}to feed\textquotedblright{} (Lewis \& Short), the source of \emph{nourrir}.
\end{cousinweb}

\begin{grammarlens}[title={What you've built}]
One more word for telling, explaining and describing at B1.
\end{grammarlens}

\subsection*{Your turn}
\begin{itemize}
\raggedright
  \item \emph{Say it:} \emph{le nourrisson}
  \item \emph{Say it:} \emph{le nourrisson}, once more
  \item \emph{Say it:} a sentence of your own with \emph{le nourrisson}
  \item \emph{Recall:} say the French for pregnancy, then the French for to give birth, then say \emph{le nourrisson} again
\end{itemize}

\begin{tcolorbox}[breakable,colback=teal!4,colframe=teal!35!black,title={Before you move on}]
What does le nourrisson mean? (\textquotedblleft{}Infant, small baby\textquotedblright{}.) Say it once more.
\end{tcolorbox}
\section[le biberon]{le biberon — baby's bottle}

\label{lesson:FR-C493-biberon}



\begin{quote}
Before the new one: say the French for to give birth, then the French for infant.
\end{quote}

\subsection*{You'll want to know: le biberon}
\textbf{le biberon} — \textquotedblleft{}baby's bottle\textquotedblright{}.

\begin{quote}
Le père donne le biberon au bébé.
\end{quote}

\begin{quote}
\emph{The father gives the baby its bottle.}
\end{quote}

\emph{Donner le biberon} — to bottle-feed.

\begin{cousinweb}
\emph{le biberon} is built on Latin \emph{bibere} \textquotedblleft{}to drink\textquotedblright{} (EtymDB; Lewis \& Short), with the ending \emph{-on}.
\end{cousinweb}

\begin{grammarlens}[title={What you've built}]
One more word for telling, explaining and describing at B1.
\end{grammarlens}

\subsection*{Your turn}
\begin{itemize}
\raggedright
  \item \emph{Say it:} \emph{le biberon}
  \item \emph{Say it:} \emph{le biberon}, once more
  \item \emph{Say it:} a sentence of your own with \emph{le biberon}
  \item \emph{Recall:} say the French for to give birth, then the French for infant, then say \emph{le biberon} again
\end{itemize}

\begin{tcolorbox}[breakable,colback=teal!4,colframe=teal!35!black,title={Before you move on}]
What does le biberon mean? (\textquotedblleft{}Baby's bottle\textquotedblright{}.) Say it once more.
\end{tcolorbox}
\section[bercer]{bercer — to rock (a baby)}

\label{lesson:FR-C493-bercer}



\begin{quote}
Before the new one: say the French for infant, then the French for baby's bottle.
\end{quote}

\subsection*{You'll want to know: bercer}
\textbf{bercer} — \textquotedblleft{}to rock (a baby)\textquotedblright{}.

\begin{quote}
Elle berce le bébé pour l'endormir.
\end{quote}

\begin{quote}
\emph{She rocks the baby to get it to sleep.}
\end{quote}

Spelling: \emph{nous berçons}. \emph{Une berceuse} — a lullaby.

\begin{cousinweb}
\emph{bercer} comes from Old French \emph{bercier} \textquotedblleft{}to rock\textquotedblright{} (EtymDB), built on \emph{bers} \textquotedblleft{}cradle\textquotedblright{}, which is usually traced to a Gaulish word.
\end{cousinweb}

\begin{grammarlens}[title={What you've built}]
One more word for telling, explaining and describing at B1.
\end{grammarlens}

\subsection*{Your turn}
\begin{itemize}
\raggedright
  \item \emph{Say it:} \emph{bercer}
  \item \emph{Say it:} \emph{bercer}, once more
  \item \emph{Say it:} a sentence of your own with \emph{bercer}
  \item \emph{Recall:} say the French for infant, then the French for baby's bottle, then say \emph{bercer} again
\end{itemize}

\begin{tcolorbox}[breakable,colback=teal!4,colframe=teal!35!black,title={Before you move on}]
What does bercer mean? (\textquotedblleft{}To rock (a baby)\textquotedblright{}.) Say it once more.
\end{tcolorbox}
